\section[Known module summand removal]{Persistent removal of a known module summand}
\label{sec:op3-module-curve-removal}

\paragraph{Status and algebraic input promise.}
This \statusproved{} draft supplies a missing curve operation; it does not
yet pay for a sequence of graph admissions. Work is in exact-real words,
using the immutable affine AVL primitives already implemented in this note.
No new external data-structure theorem is imported. A response is represented
by its genuine breakpoints and affine end rays. Consider convex,
nondecreasing, continuous piecewise-affine functions $A,B,C$ that vanish
on their left rays and satisfy $A=B+C$ in the \emph{same horizontal field}.
Write $k_A,k_B$ for the numbers of genuine breakpoints of $A,B$.

\begin{lemma}[Known convex summand removal; \statusproved]
\label{lem:op3-module-summand-removal}
Given persistent versions of $A$ and $B$ and the promise above, construct
$C=A-B$ in
\begin{equation}
 O\bigl((1+k_B)\log(2+k_A)\bigr)
 \label{eq:op3-module-summand-removal}
\end{equation}
work and allocated words. The algorithm streams only $B$ and makes
logarithmic searches and updates in $A$; it never scans all of $A$.
It deletes every breakpoint whose slope jump becomes zero and preserves
all old curve versions. In particular, repeated addition and removal of
known components need not accumulate redundant knots in the current curve.
This last assertion bounds current representation size, not the total
work or allocated history of an unbounded operation sequence.
\end{lemma}
\begin{proof}
A convex function with zero left ray has the unique hinge expansion
\[
 B(x)=\sum_{s\in\mathcal K_B}c_s(x-s)_+,\qquad c_s>0.
\]
Because $A=B+C$ and $C$ is convex, $A$ has a breakpoint at each $s$ and
its slope jump there is at least $c_s$. Removing any subset of these
hinges leaves $C$ plus the remaining positive hinges of $B$, still convex
and nondecreasing with a zero left ray. Thus no unchecked nonconvex
intermediate state is required.

Stream successive points of $B$, recovering $c_s$ from adjacent slopes
and the two end rays. In the current $A$, locate $s$ and its strict
predecessor and successor with read-only searches carrying lazy affine
tags. They give its exact left and right slopes, using stored end slopes
when a neighbor is absent. Verify that their jump is at least $c_s$.
Apply the signed suffix shear
\[
 (x,y)\longmapsto (x,y-c_s(x-s))\quad\text{for }x\geq s.
\]
This changes only the suffix boundary search path and lazy tags on
whole subtrees. Horizontal order is unchanged. The promise and the
local jump check ensure convexity; this is not an unconditional permission
to subtract an arbitrary response.

If the old jump equals $c_s$, remove the now-redundant point at $s$.
Persistent AVL deletion pushes tags and copies only search paths. With
two children, replace the point by its right subtree's first point and
rebalance on the copied path. Heights and sizes are structural and do
not change under affine tags. Each level uses a constant number of
copies and rotations, so deletion, suffix shear and strict-neighbor
queries each cost $O(\log(2+k_A))$. Every intermediate curve has at most
$k_A$ knots. Streaming and all copied nodes satisfy
\cref{eq:op3-module-summand-removal}.

If the last knot is removed, the zero left ray identifies the resulting
empty curve as identically zero. Otherwise every remaining knot has a
positive slope jump. Immutability preserves the original $A$, $B$ and
any earlier versions. The operation creates no new breakpoint position.
\end{proof}

\begin{corollary}[Extract a child in the correct module field; \statusproved]
\label{cor:op3-module-child-extraction}
From a supplied union or join response and a saved child response, obtain
the response of the union or join of the remaining children in
$O((1+|H_j|)\log(2+|H|))$ work and allocated words, retaining original
degrees, private-leaf counts and physical loads. No all-parent-curve scan
is required. This includes a sole remaining child.
\end{corollary}
\begin{proof}
At a union subtract $F_{H_j}$ directly from $F_H$. At a join first apply
the positive shear to $F_H$, obtaining $G_H=\sum_jG_j$ in the common $z$
field, and subtract the saved $G_j$ there. Then apply the negative shear
$h=z-\gamma\sum_{\ell\ne j}G_\ell(z)$ to the remainder. Its horizontal
derivative is positive because the original sum had derivative less than
$1/\gamma$ and all child derivatives are nonnegative. This is exactly
the response of the remaining children with the original diagonal degrees,
not degrees recomputed after deleting a child. The remaining physical
load and uniform field are the same algebraic inputs as before.

Each whole shear costs one persistent root copy. The removed child has
at most $2|H_j|$ breakpoints by \cref{lem:op3-module-response}, so the
removal lemma gives the stated charge. With a sole child, the positive
and negative shears cancel its saved transform, as required. Independently
applying inverse shears to summands before summing would be a different,
generally incorrect operation; common horizontal coordinates are essential.
\end{proof}

\paragraph{Measured scope and next obstruction.}
The executable \texttt{module\_curve\_removal.py} checks complete affine
identities, strict positive jumps, AVL invariants and old-version preservation
after removal and restoration. Its exact suite includes coincident knots,
negative breakpoint positions, arbitrary removal orders, complete removal
to zero and inherited lazy shears. Supplied union/join examples extract
every child and reconstruct the remaining response in the correct field.
The full record is \texttt{MODULE\_CURVE\_REMOVAL\_AUDIT.json}.

\Cref{cor:op3-module-child-extraction} prices a \emph{known} child by its
size. Replacing a large changed child can still require streaming many
knots, even when graph recognition changes few tree records. A static
small-to-large proof charges construction merges; it cannot be reused for
arbitrary future replacements without a fresh potential or structural
charge. That missing amortization, paid local discovery and general OP3
remain \statusopen. The operation does not repair a boundary that ceases
to apply a common field after a module split.
